\section{Proofs}
\label{app:proofs}
Throughout, $\Phi$ and $\varphi$ are the standard normal CDF and density, $z_\alpha=\Phi^{-1}(1-\alpha)$, and $\|\cdot\|$ is the Euclidean norm. We use two elementary facts.

\textbf{(F1)} If $x\sim\cN(\mu,\Sigma)$ and $u\in\R^d$, then $u^\top x\sim\cN(u^\top\mu,\,u^\top\Sigma u)$.

\textbf{(F2)} (Neyman--Pearson lemma, \citealp{neyman1933,lehmann2005}.) For two distributions with densities $p_0,p_1$, a test that rejects when $p_1/p_0$ exceeds a constant chosen to give size $\alpha$ is most powerful among all tests of size $\le\alpha$. Ties occur with probability $0$ when the likelihood ratio has a continuous distribution.

\subsection{Proof of \texorpdfstring{\cref{thm:np}}{Theorem 1}}
\textbf{(i)} Under \cref{ass:gauss} the densities are $p_h(x)=(2\pi)^{-d/2}|\Sigma|^{-1/2}\exp(-\frac12(x-h)^\top\Sigma^{-1}(x-h))$. Hence
\[
\log\frac{p_g(x)}{p_0(x)}=-\tfrac12(x-g)^\top\Sigma^{-1}(x-g)+\tfrac12x^\top\Sigma^{-1}x=g^\top\Sigma^{-1}x-\tfrac12g^\top\Sigma^{-1}g .
\]
By Bayes' rule the log posterior odds equal the log-likelihood ratio plus $\log\frac{\pi}{1-\pi}$, which gives $\ell_\pi$. If $\Sigma=\mathrm{blockdiag}(\Sigma_{11},\dots,\Sigma_{CC})$ over channels, then $g^\top\Sigma^{-1}x=\sum_{j}g_j^\top\Sigma_{jj}^{-1}x_j$, where $x_j,g_j\in\R^T$ are the channel-$j$ blocks. This is a sum of per-channel terms, and $g_j=0$ for $j\notin S$.

\textbf{(ii)} By (i) the likelihood ratio is a strictly increasing function of $u(x)=g^\top\Sigma^{-1}x$. By (F2), the most powerful size-$\alpha$ test rejects when $u(x)>t$, with $t$ chosen so that $P_0(u>t)=\alpha$. By (F1), under $P_0$, $u\sim\cN(0,g^\top\Sigma^{-1}\Sigma\Sigma^{-1}g)=\cN(0,D^2)$, so $t=z_\alpha D$. Neither the test nor its threshold involves $\pi$.

\textbf{(iii)} Let $P_\pi=(1-\pi)P_0+\pi P_g$ and let $\hat t_n$ be the $(1-\alpha)$-quantile of $\hat f_n$ under $P_0$. Let $t^*$ be the $(1-\alpha)$-quantile of $\ell_\pi$ under $P_0$. Since $\ell_\pi$ is an increasing affine function of $u$, $\{\ell_\pi>t^*\}=\{\phi^*=1\}$ by (ii).

\emph{Step 1 (convergence under $P_0$).} $P_0(A)\le P_\pi(A)/(1-\pi)$ for every event $A$, so $\hat f_n\to\ell_\pi$ in $P_\pi$-probability implies convergence in $P_0$-probability.

\emph{Step 2 (threshold convergence).} Let $F_n$ and $F$ be the CDFs of $\hat f_n$ and $\ell_\pi$ under $P_0$. Convergence in probability implies convergence in distribution, so $F_n(t)\to F(t)$ at every $t$, because $F$ is continuous: $\ell_\pi$ is Gaussian under $P_0$ with positive variance $D^2$. $F$ is also strictly increasing. Fix $\epsilon>0$. Then $F(t^*-\epsilon)<1-\alpha<F(t^*+\epsilon)$, and for $n$ large $F_n(t^*-\epsilon)<1-\alpha<F_n(t^*+\epsilon)$. Hence $\hat t_n\in[t^*-\epsilon,t^*+\epsilon]$, where $\hat t_n=\inf\{t:F_n(t)\ge1-\alpha\}$. So $\hat t_n\to t^*$. If $\hat t_n$ is instead the empirical quantile of $m_n\to\infty$ normal scores drawn independently of training, then conditionally on $\hat f_n$ these scores are i.i.d.\ with CDF $F_n$, and the Dvoretzky--Kiefer--Wolfowitz inequality \citep{massart1990} gives $\Pr(\sup_t|\hat F_{m_n}-F_n|>\epsilon')\le2e^{-2m_n\epsilon'^2}$ uniformly in $\hat f_n$. Taking $\epsilon'$ below the two gaps $|F_n(t^*\pm\epsilon)-(1-\alpha)|$, which stay bounded away from $0$ for large $n$, the same sandwich holds for $\hat F_{m_n}$ with probability $\to1$, so $\hat t_n\to t^*$ in probability.

\emph{Step 3 (transfer to $Q$).} $P_\pi$ has an everywhere positive density, so Lebesgue measure is absolutely continuous w.r.t.\ $P_\pi$. Hence $Q\ll P_\pi$. For finite measures this means: for every $\eta>0$ there is $\gamma>0$ with $P_\pi(A)<\gamma\Rightarrow Q(A)<\eta$. Therefore $\hat f_n\to\ell_\pi$ in $Q$-probability.

\emph{Step 4 (disagreement).} For any $\epsilon>0$, the event $\{\hat\phi_n\ne\phi^*\}$ is contained in
\[
\{|\ell_\pi-t^*|\le2\epsilon\}\ \cup\ \{|\hat f_n-\ell_\pi|>\epsilon\}\ \cup\ \{|\hat t_n-t^*|>\epsilon\}.
\]
To see this, suppose $|\ell_\pi-t^*|>2\epsilon$, $|\hat f_n-\ell_\pi|\le\epsilon$ and $|\hat t_n-t^*|\le\epsilon$. Then $\hat f_n-\hat t_n$ and $\ell_\pi-t^*$ have the same sign, so the two tests agree. The third event has probability $\to0$ (Step 2; it is eventually empty for the population quantile). The second has $Q$-probability $\to0$ (Step 3). The first has $Q$-probability $\to0$ as $\epsilon\to0$, because $\ell_\pi$ is a non-constant affine function of $x$ and so its level sets have Lebesgue measure zero, hence $Q$-measure zero. Taking $n\to\infty$ and then $\epsilon\to0$ gives $Q(\hat\phi_n\ne\phi^*)\to0$. \qed

\begin{remark}
Logistic regression on $(x,1)$ is well specified by (i), and its maximum-likelihood estimator is consistent, so it satisfies the premise of (iii). The same conclusion holds for exponential-family noise whose sufficient statistics are channel-separable: the log-likelihood ratio is then a sum of per-channel terms.
\end{remark}

\subsection{Proof of \texorpdfstring{\cref{thm:law}}{Theorem 2} and \texorpdfstring{\cref{cor:paradox}}{Corollary 3}}
Under $\cN(h,\Sigma)$, (F1) gives $u=g^\top\Sigma^{-1}x\sim\cN(g^\top\Sigma^{-1}h,\,D^2)=\cN(D^2r(h),D^2)$. Thus $\Pr_h(u>z_\alpha D)=\Pr(Z>z_\alpha-Dr(h))=\Phi(Dr(h)-z_\alpha)$ for $Z\sim\cN(0,1)$.

\emph{Special cases.} Assume independent channels with equal per-channel detectability $D_j^2=D^2/k$ for $j\in S$. For $h=g\odot\1_A$, block-diagonality gives $g^\top\Sigma^{-1}h=\sum_{j\in A}D_j^2=|A|D^2/k$, so $r=m/k$ when $|A|=m$. Shifting only channel $c$, by $c'$ times its target shift ($h=c'g_c$), gives $r=c'D_c^2/D^2=c'/k$. For $h=-g$, $r=-1$. For $h$ supported on $S^c$, $g^\top\Sigma^{-1}h=0$ because $g$ vanishes off $S$ and $\Sigma^{-1}$ is block-diagonal, so the rate is $\Phi(-z_\alpha)=\alpha$.

\emph{\Cref{cor:paradox}.} $\beta=\Phi(D-z_\alpha)$ gives $D=\Phi^{-1}(\beta)+z_\alpha$. Substituting this into the law gives the stated formula. For $r\in(0,1)$ the argument $r(\Phi^{-1}(\beta)+z_\alpha)-z_\alpha$ is strictly increasing in $\beta$ and tends to $+\infty$ as $\beta\to1$. $\Phi$ is strictly increasing. Replacing $(g,h,\Sigma)$ by $(cg,ch,\Sigma)$ with $c>1$, or by $(g,h,c\Sigma)$ with $c<1$, leaves $r(h)=g^\top\Sigma^{-1}h/g^\top\Sigma^{-1}g$ unchanged and multiplies $D$ by $c$ or $c^{-1/2}$, so $\beta$ and every look-alike rate with $r(h)\in(0,1)$ increase. \qed

\subsection{Shift statistic and proof of \texorpdfstring{\cref{thm:composite}}{Theorem 4}}
\begin{lemma}[Shift statistic]\label{lem:shift}
Let $B,I\subset[T]$ be disjoint baseline and post-onset index sets and $c=\1_I/|I|-\1_B/|B|\in\R^T$. Let $a_j=s_j\,c^\top x_j$. Under \cref{ass:gauss} with block-diagonal $\Sigma$ and a pattern $h$ whose channel-$j$ block equals $s_j\mu_j\1_I$ (and $0$ on $B$), the vector $a=(a_j)_{j\in S}$ is $\cN(\mu,\mathrm{diag}(\sigma_j^2))$ with $\sigma_j^2=c^\top\Sigma_{jj}c$.
\end{lemma}
\begin{proof}
$a_j=s_jc^\top x_j$ is linear in $x_j$, with mean $s_jc^\top(s_j\mu_j\1_I)=\mu_j$ because $c^\top\1_I=1$ and $s_j^2=1$. Its variance is $c^\top\Sigma_{jj}c$ (F1). The $a_j$ are jointly Gaussian and depend on different blocks, so block-diagonality makes them independent.
\end{proof}

\textbf{(a)} Fix $j\in S$ and let $\mu^{(j)}=\mu-\mu_je_j$, i.e.\ the near-miss without sensor $j$, with $\mu=(\delta_i)_{i\in S}$. Since $M_i\ge\delta_i$, $\cN(\mu^{(j)},\cdot)$ belongs to $\mathcal N(M)$. A deployment-valid $\phi$ is therefore in particular a level-$\alpha$ test of $\cN(\mu^{(j)},\Lambda)$ against $\cN(\mu,\Lambda)$, with $\Lambda=\mathrm{diag}(\sigma_i^2)$. These two laws differ only in the mean of coordinate $j$. By (i) of \cref{thm:np} applied to $g=\mu-\mu^{(j)}=\delta_je_j$, the log-likelihood ratio is $\delta_j(a_j-\mu^{(j)}_j)/\sigma_j^2$ plus a constant. By (F2) the most powerful level-$\alpha$ test rejects when $a_j>z_\alpha\sigma_j$, with power $\Phi(\delta_j/\sigma_j-z_\alpha)$. So no deployment-valid test exceeds this power. Minimize over $j$. \emph{Full windows:} the same argument applies to any measurable function of the whole window $x$. The target pattern $g$ and the near-miss pattern $g-g_j$ differ only in the channel-$j$ block, so with block-diagonal $\Sigma$ the likelihood ratio depends on $x_j$ only, and the most powerful test has power $\Phi(D_j-z_\alpha)$ with $D_j^2=g_j^\top\Sigma_{jj}^{-1}g_j$.

\textbf{(b)} \emph{Validity.} Let $Q\in\mathcal N(\infty)$ have support $A\subsetneq S$, and pick $j\in S\setminus A$, so $\mu_j=0$ under $Q$. Then $Q(\phi_{IU}=1)\le Q(a_j/\sigma_j>z_\alpha)=\alpha$. Under normal ($\mu=0$) the same bound holds. \emph{Power.} Under the target the $a_j$ are independent with $\Pr(a_j/\sigma_j>z_\alpha)=\Phi(D_j-z_\alpha)$, so the power is $\prod_j\Phi(D_j-z_\alpha)$. \emph{Approximation.} Write $q_j=\Phi(z_\alpha-D_j)$ for the per-coordinate miss. The IU miss is $1-\prod_j(1-q_j)\le\sum_jq_j\le k\max_jq_j$, using the union bound. By (a), every valid detector misses with probability at least $1-\min_j\Phi(D_j-z_\alpha)=\max_jq_j$.

\textbf{(c)} Let $w\ge0$, $w\ne0$, with equal $\delta,\sigma$ and $M_i=\delta$. For each $i$, the near-miss $\mu^{(i)}=\delta(\1-e_i)$ is in $\mathcal N(\delta)$, and $w^\top a\sim\cN(\delta(\|w\|_1-w_i),\sigma^2\|w\|^2)$ under it. Validity requires $\Pr(w^\top a>t)\le\alpha$, i.e.\ $t\ge\delta(\|w\|_1-w_i)+z_\alpha\sigma\|w\|$ for all $i$. The binding constraint is $i^*=\arg\min_iw_i$. The power is $\Phi\big((\delta\|w\|_1-t)/(\sigma\|w\|)\big)\le\Phi\big(\delta\min_iw_i/(\sigma\|w\|)-z_\alpha\big)$. Finally, $k\min_iw_i^2\le\|w\|^2$ gives $\min_iw_i\le\|w\|/\sqrt k$, with equality iff all $w_i$ are equal. Other members of $\mathcal N(\delta)$ have coordinatewise smaller means, hence smaller false-alarm rates (since $w\ge0$), so they do not tighten the constraint.

\textbf{(d)} Let $\psi_j$ be nondecreasing, real-valued and unbounded above. Since $k\ge2$, the vector $Ae_j$ (all sensors normal except $j$, shifted by $A>0$) has support $\{j\}\subsetneq S$, so it belongs to $\mathcal N(\infty)$ for every $A$. Fix any $t\in\R$ and $\epsilon>0$. The random variable $Y=\sum_{i\ne j}\psi_i(a_i)$ has a fixed distribution under these alternatives and is a.s.\ finite, so there is $y_\epsilon$ with $\Pr(Y<-y_\epsilon)<\epsilon$. Since $\psi_j$ is unbounded and nondecreasing, there is $u_\epsilon$ with $\psi_j(u)>t+y_\epsilon$ for all $u\ge u_\epsilon$. Under $Ae_j$, $\Pr(a_j\ge u_\epsilon)\to1$ as $A\to\infty$. Hence $\liminf_{A\to\infty}\Pr(\sum_i\psi_i(a_i)>t)\ge1-\epsilon$. So no finite threshold is valid, and the only valid choice $t=+\infty$ (never alarm) has power $0$. Linear detectors with $w_j>0$ are the case $\psi_j(u)=w_ju$. For a channel-independent model $f(X)=\sum_jf_j(X_{\cdot j})$, apply the same argument with the null member that shifts channel $j$ by $A\,s_j\1_I$. ``Unbounded along fault-direction shifts'' means $f_j(X_{\cdot j}+As_j\1_I)\to\infty$ in probability as $A\to\infty$.

\textbf{(e)} With $\psi_j=\1\{a_j>z_\alpha\sigma_j\}$ and $t=k-\frac12$, $\sum_j\psi_j>t$ iff all indicators equal $1$, i.e.\ $\phi_{IU}=1$.

\textbf{(f)} Apply \cref{thm:law} with $r=(k-1)/k$ and $D=\sqrt k\,\delta/\sigma$: $Dr=(k-1)\delta/(\sigma\sqrt k)>0$, so the rate exceeds $\Phi(-z_\alpha)=\alpha$. \qed

\subsection{Norms, Lipschitz constants, and proof of \texorpdfstring{\cref{thm:sep}}{Theorem 5}}
\begin{lemma}[Norms bound the Lipschitz constant]\label{lem:lip}
For $f(x)=\sum_iv_i\sigma(w_i^\top x+b_i)+c$ with $\sigma$ 1-Lipschitz, $|f(x)-f(y)|\le R\|x-y\|$ with $R=\sum_i|v_i|\|w_i\|$. Moreover, $\min$ over the per-unit rescalings $(v_i,w_i,b_i)\mapsto(v_i/s,sw_i,sb_i)$, $s>0$, of $\frac12\sum_i(v_i^2+\|w_i\|^2+b_i^2)$ is at least $R$. For ReLU these rescalings leave $f$ unchanged, so $\ell_2$ weight decay upper-bounds $\lambda R$ at any minimizer. For depth-$L$ networks, $\Lip(f)\le\prod_l\|W_l\|_{\mathrm{op}}\le(\|\theta\|^2/L)^{L/2}$.
\end{lemma}
\begin{proof}
$|f(x)-f(y)|\le\sum_i|v_i||\sigma(w_i^\top x+b_i)-\sigma(w_i^\top y+b_i)|\le\sum_i|v_i||w_i^\top(x-y)|\le R\|x-y\|$ by Cauchy--Schwarz. For one unit, $\min_{s>0}\frac12(v^2/s^2+s^2(\|w\|^2+b^2))=|v|\sqrt{\|w\|^2+b^2}\ge|v|\|w\|$, by AM--GM with equality at $s^2=|v|/\sqrt{\|w\|^2+b^2}$. ReLU is positively homogeneous, so the rescaling preserves $f$. For depth $L$, the Lipschitz constant of a composition is at most the product of the factors' constants. Then $\|W_l\|_{\mathrm{op}}\le\|W_l\|_F$, and AM--GM gives $\prod_l\|W_l\|_F\le(\sum_l\|W_l\|_F^2/L)^{L/2}$.
\end{proof}

\textbf{\Cref{thm:sep}(a).} By definition of the Lipschitz constant, $f(x)\ge f(x_+)-\Lip(f)\|x-x_+\|=f(x_+)-\kappa\frac{G}{\Delta}\rho\Delta=f(x_+)-\kappa\rho G$. Subtracting $f(x_0)$ and dividing by $G$ gives $p(x)\ge1-\kappa\rho$. \emph{Distances.} Take $x_0=0$ (baseline-corrected normal) and $x_+=g$, with every signature channel shifted by $\delta$ on a block $\mathcal B$ of $L$ steps. Then $\Delta=\delta\sqrt{kL}$. The channel-$j$ near-miss $x^{(j)}=g-g_j$ satisfies $\|x^{(j)}-x_+\|=\|g_j\|=\delta\sqrt L$, so $\rho=1/\sqrt k$. Displacing channel $j$'s block by $\Delta t\le L$ steps changes $\Delta t$ entries from $\delta$ to $0$ and $\Delta t$ entries from $0$ to $\delta$, so the distance is $\delta\sqrt{2\Delta t}$ and $\rho=\sqrt{2\Delta t/(kL)}$. \emph{Interpretation of $\kappa$.} Every $f$ with $f(x_+)-f(x_0)=G$ has $\Lip(f)\ge G/\Delta$, and the linear score $f(x)=\frac{G}{\Delta^2}(x_+-x_0)^\top(x-x_0)$ attains it. So $\kappa\ge1$ measures excess slope. In experiments we upper-bound $\kappa$ by $R/R_{\min}$ using \cref{lem:lip}.

\textbf{\Cref{thm:sep}(b).} For each unit and each $j$,
$|\relu(w_i^\top x_++b_i)-\relu(w_i^\top x^{(j)}+b_i)|\le|w_i^\top(x_+-x^{(j)})|=\delta|\sum_{t\in\mathcal B}w_{i,tj}|\le\delta\sum_{t\in\mathcal B}|w_{i,tj}|$, where $w_{i,tj}$ is the weight of unit $i$ on input $(t,j)$. Summing over units with weights $|v_i|$ gives $\gamma\le f(x_+)-f(x^{(j)})\le\delta\sum_i|v_i|\sum_{t\in\mathcal B}|w_{i,tj}|$. Summing over $j\in S$,
\[
k\gamma\le\delta\sum_i|v_i|\,\|w_{i,\mathcal B\times S}\|_1\le\delta\sqrt{kL}\sum_i|v_i|\,\|w_{i,\mathcal B\times S}\|_2\le\delta\sqrt{kL}\,R ,
\]
using Cauchy--Schwarz on the $kL$ coordinates of $\mathcal B\times S$. Hence $R\ge\sqrt k\gamma/(\delta\sqrt L)$. With $R_{\min}=G/\Delta=G/(\delta\sqrt{kL})$ this is $R\ge k\frac{\gamma}{G}R_{\min}$. For unequal shifts the same steps give $R\ge\gamma\sum_{j\in S}\delta_j^{-1}/\sqrt{kL}$. \qed

\begin{remark}[Exact linear limit]\label{rem:lyu}
\citet{lyu2021} show that, for linearly separable data that is symmetric (closed under $(x,y)\mapsto(-x,-y)$), gradient flow from small initialization on two-layer leaky-ReLU networks converges in direction to the linear max-margin classifier. Symmetry holds exactly after centring at the midpoint of the prototypes and symmetrizing the sample. For two prototypes $\pm(x_+-x_0)/2$ that classifier is proportional to $x_+-x_0\propto\1_S$ (in shift-statistic space), so channel near-misses sit at $p=(k-1)/k$. The midpoint threshold is $p=\frac12$, which they exceed for every $k\ge3$. We do not rely on this result for non-symmetric data. \Cref{thm:sep} holds for any network, in terms of its measured $\kappa$.
\end{remark}

\begin{proposition}[Necessary near-miss fraction under weight decay]\label{prop:pstar}
In shift-statistic space, take noise-free prototypes with mass $q$ at $x_0=0$ (label $-1$), $q$ at $x_+=\delta\1_S$ (label $+1$), and $p/k$ at each channel near-miss (label $-1$), with $2q+p=1$. Let $\ell(z)=\log(1+e^{-z})$ and consider two-layer ReLU networks with objective $J(\theta)=\sum\mathrm{mass}\cdot\ell(yf_\theta(x))+\lambda R(\theta)$. Define
\begin{align*}
J_{\mathrm{cheap}}(p)&=\min_{m_0,m_1\ge0}\ q\ell(m_0)+q\ell(m_1)+p\,\ell\big(-m_1+\tfrac{m_0+m_1}{k}\big)+\lambda\tfrac{m_0+m_1}{\delta\sqrt k},\\
B(p)&=\inf_{m,m_0\in\R}\ q\ell(m)+q\ell(m_0)+\lambda\max\Big(\tfrac{\sqrt k\,m_+}{\delta},\tfrac{|m+m_0|}{\delta\sqrt k}\Big).
\end{align*}
If $J_{\mathrm{cheap}}(p)<B(p)$, every minimizer of $J$ flags at least one near-miss ($f(x^{(j)})>0$). The same holds with $\lambda R$ replaced by $\ell_2$ weight decay $\frac{\lambda}{2}\sum_i(v_i^2+\|w_i\|^2+b_i^2)$ with an unpenalized output bias: by \cref{lem:lip} the weight decay of every $\theta$ is at least $\lambda R(\theta)$, and for the one-unit network of the upper bound, which has hidden bias $0$, the best rescaling makes it equal to $\lambda R$.
\end{proposition}
\begin{proof}
\emph{Upper bound.} The one-unit network $f(x)=v\,\relu(\1_S^\top x/\sqrt k)-m_0$ with $v=(m_0+m_1)/(\delta\sqrt k)$ has $f(x_0)=-m_0$, $f(x_+)=m_1$ and $f(x^{(j)})=v\delta(k-1)/\sqrt k-m_0=m_1-(m_0+m_1)/k$. Its path norm is $v$, since $\|\1_S/\sqrt k\|=1$. So $\min J\le J_{\mathrm{cheap}}(p)$. \emph{Lower bound.} Let $f$ reject every near-miss, $f(x^{(j)})\le0$, and write $m=f(x_+)$, $m_0=-f(x_0)$. The near-miss losses are nonnegative, so they can be dropped. \Cref{thm:sep}(b) with $\gamma=m_+$ ($L=1$) gives $R\ge\sqrt k\,m_+/\delta$. \Cref{lem:lip} with $\|x_+-x_0\|=\delta\sqrt k$ gives $R\ge|m+m_0|/(\delta\sqrt k)$. Hence $J(\theta)\ge B(p)$, and when $J_{\mathrm{cheap}}(p)<B(p)$ no near-miss-rejecting network attains $\min J$.
\end{proof}
Numerically ($\delta=1$), the largest $p$ with $J_{\mathrm{cheap}}<B$ is $p^*=2\times10^{-4},1.3\times10^{-3},8.8\times10^{-3},1.7\times10^{-2}$ for $\lambda=10^{-4},10^{-3},10^{-2},3\times10^{-2}$ at $k=3$, and $p^*$ grows with $k$ (to $6.1\times10^{-2}$ at $k=8$, $\lambda=3\times10^{-2}$). This is a \emph{necessary} condition only. With noisy data, gradient-trained networks needed far more near-misses (\cref{tab:recipe}).

\subsection{Proof of \texorpdfstring{\cref{thm:teach}}{Theorem 6}}
\textbf{Setting.} The labelled examples are configurations labelled by the target rule $h^*$. A labelled set $\mathcal L$ identifies $h^*$ within $\mathcal H$ if $h^*$ is the only $h\in\mathcal H$ with $h(c)=h^*(c)$ for all $c\in\mathcal L$. \emph{Channel axis.} Configurations are $A\subseteq S$ (sensors shifted, simultaneously). $\mathcal H_{\mathrm{ch}}$ is the set of $h:2^S\to\{0,1\}$ that are monotone ($A\subseteq A'\Rightarrow h(A)\le h(A')$) with $h(S)=1$. The target is $h^*_{\mathrm{ch}}(A)=\1\{A=S\}$, and training data label $A=\emptyset$ (normal) as $0$. \emph{Time axis.} All of $S$ is shifted on nonempty half-open intervals $I_j=[l_j,u_j)$. Configurations are the realizable overlap graphs $\mathcal G$ on vertex set $S$. $\mathcal H_{\mathrm t}$ is the set of $h$ monotone in the edge set with $h(K_k)=1$. The target ``some moment at which all $k$ are shifted'' is $h^*_{\mathrm t}$.

\textbf{(a)} \emph{Sufficiency.} Let $\mathcal L\ni\emptyset$ contain every $S\setminus\{j\}$. Any $A\ne S$ omits some $j$, so $A\subseteq S\setminus\{j\}$, and monotonicity gives $h(A)\le h(S\setminus\{j\})=0$ for every consistent $h$. Together with $h(S)=1$ this forces $h=h^*_{\mathrm{ch}}$. \emph{Necessity.} If $S\setminus\{i\}\notin\mathcal L$, let $h_i(A)=\1\{A\supseteq S\setminus\{i\}\}$. It is monotone, $h_i(S)=1$, and $h_i(\emptyset)=0$ since $k\ge2$. It differs from $h^*_{\mathrm{ch}}$ exactly at $A=S\setminus\{i\}$, so it agrees with $h^*_{\mathrm{ch}}$ on $\mathcal L$ and $\mathcal L$ does not identify. (This characterization of teaching sets for monotone rules is classical, cf.\ \citealp{goldman1995}.)

\textbf{(b)} \emph{Helly on the line.} If $I_i\cap I_j\ne\emptyset$ for all $i,j$, then $l_i<u_j$ for all $i,j$, so $\max_il_i<\min_ju_j$, and every point of $[\max_il_i,\min_ju_j)$ lies in all intervals. Conversely a common point lies in every pairwise intersection. Hence $h^*_{\mathrm t}(\mathcal G)=\1\{\mathcal G=K_k\}$, the conjunction of the $\binom k2$ edge indicators. \emph{Sufficiency.} If $\mathcal L$ contains $K_k-e$ for every edge $e$, any realizable $\mathcal G\ne K_k$ misses some edge $e$, so $\mathcal G\subseteq K_k-e$ and $h(\mathcal G)\le h(K_k-e)=0$. \emph{Necessity.} If $K_k-e\notin\mathcal L$, let $h_e(\mathcal G)=\1\{\mathcal G\supseteq K_k-e\}$. It is monotone with $h_e(K_k)=1$ and differs from $h^*_{\mathrm t}$ only at $K_k-e$, so it is consistent with $\mathcal L$. \emph{Realizability with equal durations.} For $e=\{a,b\}$ take $I_a=[0,1)$, $I_b=[1,2)$ and $I_j=[\frac12,\frac32)$ otherwise. Then $I_a\cap I_j=[\frac12,1)$, $I_b\cap I_j=[1,\frac32)$, $I_j\cap I_{j'}=I_j$ and $I_a\cap I_b=\emptyset$, so $\mathcal G=K_k-e$. (\Cref{alg:cis} uses exactly this construction scaled to the post-onset window.)

\textbf{(c)} The fully staggered configuration has $\mathcal G=\emptyset$, the minimum of the edge order. Labelling it $0$ excludes exactly the rules with $h(\emptyset)=1$. By monotonicity that is only the constant rule $h\equiv1$. For $k\ge3$, $K_k-e\ne\emptyset$, so every witness $h_e$ in the necessity proof of (b) satisfies $h_e(\emptyset)=0$ and survives. Fully staggered negatives therefore neither identify the rule nor replace any pair-staggered pattern. (For $k=2$, $K_2-e=\emptyset$ and the two notions coincide.)

\textbf{(d)} Let $\mathcal C_{\mathrm{ch}}$ and $\mathcal C_{\mathrm t}$ be the two configuration sets, sharing only the target (all $k$ sensors shifted and overlapping). Let $\mathcal H$ contain the rules on $\mathcal C_{\mathrm{ch}}\cup\mathcal C_{\mathrm t}$ whose restrictions lie in $\mathcal H_{\mathrm{ch}}$ and $\mathcal H_{\mathrm t}$. The constraints decouple, since every member equals $1$ at the shared target. So $\mathcal L$ identifies $h^*$ in $\mathcal H$ iff $\mathcal L\cap\mathcal C_{\mathrm{ch}}$ and $\mathcal L\cap\mathcal C_{\mathrm t}$ identify on their axes. By (a) and (b) this needs and is achieved by the $k$ channel and the $\binom k2$ pair-staggered near-misses. \qed

\begin{remark}[Joint monotonicity]\label{rem:joint}
Suppose instead that rules are monotone in the joint order $(A,\mathcal G)\preceq(A',\mathcal G')$ iff $A\subseteq A'$ and $\mathcal G\subseteq\mathcal G'$, over all configurations. Then the channel near-miss $(S\setminus\{j\},K_{S\setminus\{j\}})$ lies below $(S,K_k-\{i,j\})$ for any $i\ne j$. So the $\binom k2$ pair-staggered negatives would already force every channel near-miss to $0$. We do not assume joint monotonicity for learned networks, and we did not test pair-staggered negatives alone. Our real-data remedy (CIS) supplies both types.
\end{remark}

\subsection{Proof of \texorpdfstring{\cref{prop:samples}}{Proposition 7}}
Let $\bar x_j$ and $\bar y_j$ be the sample means of sensor $j$ over $n$ i.i.d.\ near-misses of type $j$ and $n$ targets, with true means $\nu_j$ and $\nu_j+\Delta_j$ and noise $\cN(0,\sigma_j^2)$. The threshold is $\hat c_j=(\bar x_j+\bar y_j)/2$, and the midpoint is $c_j=\nu_j+\Delta_j/2$. Then $|\hat c_j-c_j|\le\frac12(|\bar x_j-\nu_j|+|\bar y_j-\nu_j-\Delta_j|)\le\max(|\bar x_j-\nu_j|,|\bar y_j-\nu_j-\Delta_j|)$. The Gaussian tail bound gives $\Pr(|\bar x_j-\nu_j|>t)\le2e^{-nt^2/(2\sigma_j^2)}$, and similarly for $\bar y_j$. With $t=\Delta_j/4$, $\Pr(|\hat c_j-c_j|>\Delta_j/4)\le4e^{-n\Delta_j^2/(32\sigma_j^2)}$. A union bound over $j\in S$ shows that all thresholds are within $\Delta_j/4$ of their midpoints with probability $\ge1-4k\,e^{-n\min_j\Delta_j^2/(32\max_j\sigma_j^2)}\ge1-\eta$ under the stated condition on $n$. On this event $\hat c_j\ge\nu_j+\Delta_j/4$. A new near-miss of type $j$ is flagged only if its sensor-$j$ value exceeds $\hat c_j$, which has probability $\le\Phi(-\Delta_j/(4\sigma_j))$. Symmetrically $\hat c_j\le\nu_j+\frac34\Delta_j$, so the target fails sensor $j$ with probability $\le\Phi(-\Delta_j/(4\sigma_j))$. A union bound over $j$ bounds the miss probability. If the $n$ values are consecutive draws of a stationary AR(1) sequence with $\phi\in[0,1)$, the variance of their mean is $\frac{\sigma_j^2}{n}\big(1+2\sum_{h=1}^{n-1}(1-\frac hn)\phi^h\big)\le\frac{\sigma_j^2}{n}\frac{1+\phi}{1-\phi}$ for every $n$ (\cref{lem:ar}). The mean is still Gaussian, since it is a linear functional, so the tail bound holds with $n$ replaced by $n(1-\phi)/(1+\phi)$. \qed

\subsection{Window statistics and temporal dependence}
\begin{lemma}[Window statistics are timing- and sign-blind]\label{lem:blind}
Let $\psi:\R^T\to\R^q$ be invariant to permutations of the time index (e.g.\ mean, standard deviation, minimum, maximum, quantiles), and $\Psi(X)=(\psi(X_{\cdot1}),\dots,\psi(X_{\cdot C}))$. Suppose the noise has exchangeable time entries within each channel and independent channels. (i) If two patterns $G,G'$ have, channel by channel, the same multiset of entries, then $\Psi(W+G)\overset{d}{=}\Psi(W+G')$. In particular, simultaneous, fully staggered and pair-staggered patterns with equal per-channel durations are indistinguishable by any detector built on $\Psi$. (ii) If $W\overset d=-W$ and $\psi(-x)=\psi(x)$ (dispersion statistics), then $\psi(W+G)\overset d=\psi(W-G)$.
\end{lemma}
\begin{proof}
(i) For each $j$ choose a permutation $\pi_j$ with $G'_{\cdot j}=G_{\pi_j(\cdot),j}$. Exchangeability gives $W_{\cdot j}\overset d=W_{\pi_j(\cdot),j}$, and independence across channels makes this joint over $j$. Hence $(W+G')_{\cdot j}\overset d=(W+G)_{\pi_j(\cdot),j}$ jointly, and $\psi$ is permutation-invariant. (ii) $\psi(W+G)=\psi(-W-G)\overset d=\psi(W-G)$.
\end{proof}
With AR noise, or with order-dependent statistics such as a slope or a last-$m$ mean, (i) holds only approximately. Consistent with this, the window-statistics models flag consecutive-segment staggered faults at $0.57$--$1.00$ on the real targets of \cref{tab:realall}.

\begin{lemma}[Temporal dependence enters through the effective noise]\label{lem:ar}
For stationary noise with autocovariance $\gamma(\cdot)$, the shift statistic of \cref{lem:shift} has $\sigma_j^2=c^\top\Gamma c$ with $\Gamma_{st}=\gamma(s-t)$. For AR(1) noise with coefficient $\phi\in(-1,1)$ and marginal variance $\sigma_x^2$, the variance of an $L$-step mean is $\sigma_x^2L^{-2}\sum_{s,t}\phi^{|s-t|}$, and $L$ times this tends to $\sigma_x^2(1+\phi)/(1-\phi)$ as $L\to\infty$. All results above hold with these effective variances. For a post-onset mean against a known baseline ($c=\1_I/|I|$) and $\phi\in[0,1)$, a larger $\phi$ raises $\sigma_j^2$ and lowers $D$, which moves a detector \emph{along} the curve of \cref{cor:paradox}: detection and look-alike alarms both fall. This is not monotone for every statistic. Correlated noise partly cancels in the baseline-corrected difference, and for a unit step on the last 28 of $T=48$ steps (unit marginal variance) the full-window $D^2$ falls from $28$ at $\phi=0$ to $6.7$ at $\phi=0.9$ and rises to $50$ at $\phi=0.99$.
\end{lemma}
\begin{proof}
$a_j$ is a linear functional of $x_j$ (F1). For AR(1), $\gamma(h)=\sigma_x^2\phi^{|h|}$. Then $L\cdot L^{-2}\sum_{s,t=1}^L\phi^{|s-t|}=1+2\sum_{h=1}^{L-1}(1-h/L)\phi^h\to1+2\phi/(1-\phi)=(1+\phi)/(1-\phi)$ by dominated convergence. For $\phi\ge0$ every term is nonnegative and at most $\phi^h$, which gives the bound used in \cref{prop:samples}, and each term is nondecreasing in $\phi$, which gives the monotonicity claim.
\end{proof}

%=====================================================================
\section{Experimental details}
\label{app:exp}
\paragraph{Synthetic experiments} (code: \texttt{paper\_synth.py}, \texttt{theory\_checks.py}). All thresholds are the $0.99$ quantile of the score on $20{,}000$ fresh normal samples, and rates are estimated on $4{,}000$--$20{,}000$ fresh samples per construct. \emph{Law/paradox (S1):} $d=20$, $k=3$, $\sigma=1$, $\delta/\sigma\in\{0.5,\dots,3\}$, $2{,}000+2{,}000$ training samples. Learners: full-batch gradient-descent logistic regression ($2{,}000$ steps, step size $0.5$), and a two-layer ReLU net (64 units, initialization scale $10^{-2}$, Adam, $1{,}500$ steps, learning rate $10^{-2}$). 5 seeds. \emph{Composite (C3):} shift statistics with $k=3$, $\sigma=1$, $4\times10^5$ draws per construct; ``unbounded'' look-alikes use shifts of $50\delta$. \emph{Separated regime (S2):} $d=10$, $\sigma=0.15$, $\delta=1$, $500+500$ samples, two-layer ReLU (64 units, initialization $10^{-3}$, Adam, $3{,}000$ steps, $10^{-2}$), with 100 near-misses per type in the ``with near-misses'' runs. 3 seeds. \emph{Temporal (S3):} raw windows $T=60$, $C=6$, $k=3$, $\sigma=0.5$, $\delta=1$, shift duration $L\sim\mathcal U\{12,\dots,18\}$, random onset; $1{,}500+1{,}500$ samples plus $600$ timing negatives; two-layer ReLU on the flattened window (64 units, Adam, $1{,}200$ steps, $3\times10^{-3}$, weight decay $10^{-4}$). 5 seeds. \emph{Recipe (S4):} $d=10$, $k=3$, $\sigma=0.3$, $1{,}000+1{,}000$ samples plus $c\in\{0,4,16,64,256\}$ near-misses per type, weight decay $\in\{0,10^{-3}\}$. 3 seeds.

\paragraph{Toy replica} (\texttt{run\_toy.py}). 20 channels, 48 steps, AR(1) with $\phi=0.7$ and unit variance, a shared slow load factor ($\phi=0.97$, weight $0.6$) on channels 0--5. The target lowers channels 0, 1, 2 by $2.5$ from a random onset in $[10,25]$. The per-window baseline (mean of the first 8 steps) is subtracted. 4,000 training windows (50/50), 2,000 normal validation windows for calibration, 2,000 test windows per construct. 3 seeds (data and initialization).

\paragraph{Real data} (\texttt{real\_data.py}, \texttt{run\_real.py}). LBNL fan-coil and single-duct AHU simulations \citep{granderson2023}, each a set of one-year runs: one fault-free run and one run per fault and severity, with identical weather and schedule. Preprocessing: resampling to 15 minutes. We drop setpoint channels, channels constant in the fault-free run, and channels that shift in the same direction in $\ge80\%$ of fault families (a fault-free/faulty configuration mismatch in the AHU benchmark). Each channel is standardized by its fault-free statistics. \emph{Signature}: the $k=3$ channels with the largest mean standardized shift while the fault is visibly present, skipping near-duplicates ($|\rho|>0.98$; disabled for the AHU coil valve, where valve position and command diverge by design). A window qualifies as a training positive only if every signature channel moves by $\ge0.5$ s.d.\ ($0.3$ for the AHU coil valve) in its fault direction after onset. Windows are 48 steps (12 h), with onset uniform in $[10,25]$ and the first-8-step baseline subtracted. Train, validation and test splits use alternating 7-day blocks. Training uses 4,000 windows (50/50), with half of the normal windows time-matched to positives with the fault removed. Testing uses 1,000 windows per construct and up to 300 windows per other fault file. \emph{Clean mode} applies the fault effect only to signature channels. \emph{Real mode} keeps every channel's real deviation.

\paragraph{Models.} Window statistics: mean, min, max, s.d., slope and last-8 mean per channel. Logistic regression ($C=1$, standardized), random forest (200 trees, min leaf 2), gradient boosting (scikit-learn defaults), MLP on every 4th step (hidden $128,64$, Adam, early stopping). LSTM: 32 hidden units, BPTT, Adam ($3\times10^{-3}$), batch 64, gradient clipping 5, up to 30 epochs, early stopping on 15\% validation loss with patience 6. Channel-independent model: shared per-channel encoder $z_c=\relu(x_cW_1+b_1)$ (32 units) with a linear head $\sum_cA_c^\top z_c$. Capped CI: $\beta\sum_c\tanh(A_c^\top z_c+d_c)$, with the same encoder and a learnable $\beta$ initialized at 3. All neural models are implemented in NumPy and pass finite-difference gradient checks. \emph{Reference rule}: each signature channel minus a ridge prediction from the non-signature channels (fit on normal windows), scaled by its residual s.d., and $\max_\tau\min_{j\in S}$ of the 8-step window means after step 8. Calibrated like the models.

\paragraph{Remedies.} Each fix adds $n_{\mathrm{train}}/2=2{,}000$ label-0 windows: channel near-misses (all-but-one, cycling over the left-out channel); fully staggered (half ``consecutive segments'', half ``separate short dips''); CIS (half channel near-misses, half pair-staggered, cycling over pairs); generic negatives (the fault's effect moved onto $k$ unrelated channels). The recipe sweep adds $c\in\{0,3,10,30,100,300,1000\}$ near-misses per type to an otherwise identical training set. The signal sweep multiplies $D$ by a factor in both training and test.

\paragraph{Compute.} All experiments run on a 2-core CPU. The synthetic suite takes about 30 minutes, one toy seed about 10 minutes, and one real-data target 10--60 minutes.

%=====================================================================
\section{Additional results}
\label{app:more}
\begin{table}[h]\centering\small
\caption{Synthetic law check (S1; mean$\pm$s.d., 5 seeds). \emph{Law}: \cref{cor:paradox} evaluated at the model's own measured detection rate (no fitting).}\label{tab:law}
\setlength{\tabcolsep}{3pt}\begin{tabular}{llccccc}\toprule
Model & $\delta/\sigma$ & Detection & 2-of-3 & law & 1-of-3 & law\\\midrule
GD logistic & 0.5 & 0.070$\pm$0.005 & 0.040$\pm$0.003 & 0.039 & 0.021$\pm$0.002 & 0.021\\
 & 1.0 & 0.274$\pm$0.012 & 0.122$\pm$0.007 & 0.120 & 0.041$\pm$0.003 & 0.040\\
 & 1.5 & 0.606$\pm$0.015 & 0.278$\pm$0.012 & 0.276 & 0.074$\pm$0.004 & 0.072\\
 & 2.0 & 0.869$\pm$0.011 & 0.491$\pm$0.017 & 0.489 & 0.123$\pm$0.007 & 0.120\\
 & 2.5 & 0.974$\pm$0.002 & 0.706$\pm$0.013 & 0.699 & 0.190$\pm$0.009 & 0.184\\
 & 3.0 & 0.998$\pm$0.000 & 0.863$\pm$0.006 & 0.867 & 0.273$\pm$0.010 & 0.272\\
GD 2-layer ReLU & 0.5 & 0.026$\pm$0.004 & 0.019$\pm$0.002 & 0.019 & 0.014$\pm$0.001 & 0.014\\
 & 1.0 & 0.113$\pm$0.007 & 0.059$\pm$0.005 & 0.057 & 0.027$\pm$0.002 & 0.025\\
 & 1.5 & 0.379$\pm$0.026 & 0.164$\pm$0.009 & 0.163 & 0.050$\pm$0.002 & 0.049\\
 & 2.0 & 0.760$\pm$0.012 & 0.370$\pm$0.009 & 0.381 & 0.093$\pm$0.004 & 0.094\\
 & 2.5 & 0.941$\pm$0.004 & 0.582$\pm$0.022 & 0.605 & 0.141$\pm$0.009 & 0.151\\
 & 3.0 & 0.991$\pm$0.003 & 0.785$\pm$0.015 & 0.793 & 0.222$\pm$0.009 & 0.226\\
\bottomrule\end{tabular}\end{table}

\begin{table}[h]\centering\small
\caption{Composite null (C3, $k=3$, $4\times10^5$ draws). \emph{Size}: worst false-alarm rate over look-alikes, with shifts bounded by $\delta$ or unbounded ($50\delta$). Additive detectors are calibrated either on normal data or for the bounded null (\cref{thm:composite}(c)).}\label{tab:composite}
\setlength{\tabcolsep}{3pt}\begin{tabular}{ccccccccc}\toprule
 & \multicolumn{3}{c}{Additive, normal-calibrated} & \multicolumn{2}{c}{Additive, bounded-null} & \multicolumn{2}{c}{Intersection-union} & Bound\\
$\delta/\sigma$ & size ($\le\delta$) & size ($\infty$) & power & size ($\infty$) & power & size ($\infty$) & power (pred.) & (a)\\\midrule
2 & 0.493 & 1.000 & 0.872 & 1.000 & 0.121 & 0.010 & 0.051 (0.052) & 0.372\\
3 & 0.872 & 1.000 & 0.998 & 1.000 & 0.276 & 0.010 & 0.423 (0.421) & 0.750\\
4 & 0.989 & 1.000 & 1.000 & 1.000 & 0.493 & 0.010 & 0.866 (0.865) & 0.953\\
5 & 1.000 & 1.000 & 1.000 & 1.000 & 0.712 & 0.010 & 0.989 (0.989) & 0.996\\
\bottomrule\end{tabular}\end{table}

\begin{table}[h]\centering\small
\caption{Separated regime (S2; mean$\pm$s.d., 3 seeds). Left block: normal-vs-target training. \emph{Position}: $p(x^{(j)})$ of channel near-misses at the prototypes. \emph{Bound}: $1-(R/R_{\min})/\sqrt k$, a valid lower bound on the position because $\kappa\le R/R_{\min}$ (\cref{lem:lip}). \emph{Flagged}: near-miss flag rate at 1\% false alarms. Right block: retrained with 100 near-misses per type. \emph{Norm ratio}: $R/G$ after over $R/G$ before. \cref{thm:sep}(b) requires $R/R_{\min}\ge k\gamma/G$ with $\gamma$ the smallest target-minus-near-miss score gap; it held in every run. Detection after retraining is 1.000 throughout.}\label{tab:sep}
\setlength{\tabcolsep}{3pt}\begin{tabular}{cccccc|cccc}\toprule
$k$ & linear $\frac{k-1}{k}$ & position & bound & $R/R_{\min}$ & flagged & norm ratio & $R/R_{\min}$ & $k\gamma/G$ & flagged\\\midrule
2 & 0.500 & 0.565$\pm$0.049 & 0.049 & 1.34 & 0.980 & 3.00$\pm$0.34 & 4.03 & 1.89 & 0.203\\
3 & 0.667 & 0.727$\pm$0.009 & 0.267 & 1.27 & 1.000 & 3.82$\pm$0.53 & 4.85 & 2.24 & 0.600\\
4 & 0.750 & 0.808$\pm$0.023 & 0.404 & 1.19 & 1.000 & 5.96$\pm$0.84 & 7.11 & 2.77 & 0.797\\
6 & 0.833 & 0.886$\pm$0.010 & 0.505 & 1.21 & 1.000 & 7.05$\pm$1.14 & 8.56 & 3.53 & 0.978\\
8 & 0.875 & 0.911$\pm$0.004 & 0.574 & 1.20 & 1.000 & 5.81$\pm$0.78 & 7.00 & 1.25 & 1.000\\
\bottomrule\end{tabular}\end{table}

\begin{table}[h]\centering\small
\caption{Recipe curves: 2-of-3 look-alike flag rate vs.\ channel near-misses per type. Top: synthetic two-layer ReLU (S4; mean$\pm$s.d., 3 seeds). Bottom: real data, LSTM, one seed (detection $\ge0.99$ throughout).}\label{tab:recipe}
\setlength{\tabcolsep}{3pt}\begin{tabular}{lccccccc}\toprule
Synthetic, per type & 0 & 4 & 16 & 64 & 256 & & \\\midrule
weight decay 0 & 0.87$\pm$0.03 & 0.53$\pm$0.03 & 0.48$\pm$0.17 & 0.44$\pm$0.08 & 0.34$\pm$0.03 & & \\
weight decay 0.001 & 0.93$\pm$0.00 & 0.90$\pm$0.01 & 0.77$\pm$0.05 & 0.60$\pm$0.05 & 0.46$\pm$0.03 & & \\
\midrule Real, per type & 0 & 3 & 10 & 30 & 100 & 300 & 1000\\\midrule
FCU fan (clean), wd $10^{-3}$ & 0.96 & 0.89 & 0.75 & 0.57 & 0.26 & 0.06 & 0.01\\
FCU fan (clean), wd 0 & 0.95 & 0.92 & 0.82 & 0.62 & 0.54 & 0.44 & 0.08\\
FCU fan (real), wd $10^{-3}$ & 1.00 & 1.00 & 0.69 & 0.65 & 0.41 & 0.14 & 0.00\\
FCU fan (real), wd 0 & 0.99 & 1.00 & 0.78 & 0.74 & 0.55 & 0.32 & 0.12\\
FCU valve, wd $10^{-3}$ & 1.00 & 1.00 & 0.89 & 0.72 & 0.51 & 0.11 & 0.09\\
FCU valve, wd 0 & 1.00 & 1.00 & 1.00 & 0.86 & 0.81 & 0.64 & 0.06\\
AHU coil valve, wd $10^{-3}$ & 0.81 & 0.78 & 0.84 & 0.68 & 0.60 & 0.49 & 0.50\\
AHU coil valve, wd 0 & 0.74 & 0.79 & 0.83 & 0.67 & 0.53 & 0.54 & 0.47\\
\bottomrule\end{tabular}\end{table}

\begin{table}[h]\centering\small
\caption{Toy replica, all models, normal-vs-target training (mean$\pm$s.d., 3 seeds). Columns \emph{law} give \cref{cor:paradox} at each model's mean detection rate with equal shares ($r=m/k$).}\label{tab:toyall}
\setlength{\tabcolsep}{2.5pt}\resizebox{\linewidth}{!}{\begin{tabular}{lcccccccccc}\toprule
Model & Detection & 1-of-3 & law & 2-of-3 & law & BIG1 & Fully st. & Pair-st. & UP & OFF\\\midrule
LogReg (stats) & 0.98$\pm$0.00 & 0.26$\pm$0.04 & 0.19 & 0.79$\pm$0.02 & 0.73 & 1.00$\pm$0.00 & 0.51$\pm$0.07 & 0.75$\pm$0.05 & 0.00$\pm$0.00 & 0.01$\pm$0.00\\
Random forest (stats) & 0.95$\pm$0.01 & 0.10$\pm$0.03 & 0.16 & 0.58$\pm$0.12 & 0.62 & 0.14$\pm$0.07 & 0.39$\pm$0.07 & 0.58$\pm$0.07 & 0.00$\pm$0.00 & 0.01$\pm$0.01\\
Grad. boosting (stats) & 0.97$\pm$0.01 & 0.14$\pm$0.02 & 0.18 & 0.73$\pm$0.08 & 0.68 & 0.26$\pm$0.09 & 0.55$\pm$0.03 & 0.75$\pm$0.04 & 0.00$\pm$0.00 & 0.01$\pm$0.00\\
MLP & 0.94$\pm$0.01 & 0.15$\pm$0.01 & 0.15 & 0.60$\pm$0.02 & 0.61 & 0.82$\pm$0.01 & 0.16$\pm$0.01 & 0.40$\pm$0.02 & 0.00$\pm$0.00 & 0.01$\pm$0.00\\
LSTM & 0.94$\pm$0.02 & 0.13$\pm$0.00 & 0.15 & 0.55$\pm$0.03 & 0.60 & 0.65$\pm$0.08 & 0.14$\pm$0.02 & 0.27$\pm$0.02 & 0.00$\pm$0.00 & 0.00$\pm$0.00\\
Channel-indep. & 0.97$\pm$0.01 & 0.20$\pm$0.03 & 0.18 & 0.73$\pm$0.02 & 0.68 & 0.99$\pm$0.01 & 0.19$\pm$0.02 & 0.44$\pm$0.07 & 0.00$\pm$0.00 & 0.00$\pm$0.00\\
Capped CI & 0.90$\pm$0.01 & 0.07$\pm$0.01 & 0.13 & 0.42$\pm$0.05 & 0.54 & 0.09$\pm$0.01 & 0.09$\pm$0.01 & 0.23$\pm$0.05 & 0.00$\pm$0.00 & 0.01$\pm$0.00\\
Reference rule & 0.89$\pm$0.03 & 0.05$\pm$0.01 & 0.13 & 0.21$\pm$0.04 & 0.52 & 0.05$\pm$0.02 & 0.09$\pm$0.03 & 0.33$\pm$0.06 & 0.00$\pm$0.00 & 0.02$\pm$0.00\\
\bottomrule\end{tabular}}\end{table}

\begin{table}[h]\centering\small
\caption{Toy replica: 2-of-3 / pair-staggered flag rates per model and remedy (mean over 3 seeds).}\label{tab:toyremedy}
\setlength{\tabcolsep}{3pt}\begin{tabular}{lccccc}\toprule
Model & normal only & + fully stag. & + generic (OE) & + channel NM & + CIS\\\midrule
LogReg (stats) & 0.79 / 0.75 & 0.62 / 0.13 & 0.79 / 0.79 & 0.61 / 0.39 & 0.59 / 0.09\\
Random forest (stats) & 0.58 / 0.58 & 0.41 / 0.15 & 0.67 / 0.64 & 0.16 / 0.38 & 0.19 / 0.12\\
Grad. boosting (stats) & 0.73 / 0.75 & 0.46 / 0.13 & 0.77 / 0.76 & 0.18 / 0.53 & 0.21 / 0.10\\
MLP & 0.60 / 0.40 & 0.52 / 0.29 & 0.56 / 0.36 & 0.26 / 0.29 & 0.34 / 0.10\\
LSTM & 0.55 / 0.27 & 0.40 / 0.12 & 0.56 / 0.26 & 0.15 / 0.08 & 0.23 / 0.04\\
Channel-indep. & 0.73 / 0.44 & 0.65 / 0.11 & 0.73 / 0.44 & 0.33 / 0.23 & 0.39 / 0.05\\
Capped CI & 0.42 / 0.23 & 0.46 / 0.07 & 0.47 / 0.23 & 0.19 / 0.20 & 0.27 / 0.05\\
\bottomrule\end{tabular}\end{table}


\begin{table}[h]\centering\small
\caption{Real data (LBNL), normal-vs-target training, all models, one seed. Clean mode. \emph{Pairwise law}: 2-of-3 rate predicted from the detection rate and the single-sensor rates (each sensor's share $r_j$ read from its single-sensor flag rate), a test of additivity. FCU values saturate near 1; the AHU coil is informative.}\label{tab:realall}
\setlength{\tabcolsep}{2.4pt}
\resizebox{\linewidth}{!}{%
\begin{tabular}{llccccccccc}
\toprule
Target & Model & Det. & 1-of-3 & 2-of-3 & pairwise law & BIG1 & Fully st. & Pair-st. & UP & OFF\\
\midrule
FCU fan & LogReg (stats) & 1.00 & 0.81 & 1.00 & 1.00 & 0.98 & 0.95 & 0.97 & 0.00 & 0.22\\
 & Random forest & 1.00 & 1.00 & 1.00 & 1.00 & 1.00 & 1.00 & 1.00 & 0.98 & 0.33\\
 & Grad.\ boosting & 1.00 & 0.65 & 1.00 & 1.00 & 0.65 & 0.94 & 0.94 & 0.13 & 0.07\\
 & MLP & 0.99 & 0.69 & 0.95 & 1.00 & 0.96 & 0.76 & 0.87 & 0.00 & 0.17\\
 & LSTM & 1.00 & 0.79 & 0.97 & 1.00 & 0.98 & 0.85 & 0.92 & 0.48 & 0.61\\
 & Channel-indep. & 1.00 & 0.88 & 0.99 & 1.00 & 1.00 & 0.95 & 0.95 & 0.00 & 0.00\\
 & Capped CI & 1.00 & 0.96 & 1.00 & 1.00 & 1.00 & 0.99 & 0.98 & 0.41 & 0.14\\
 & Reference rule & 1.00 & 0.08 & 0.35 & 0.25 & 0.06 & 0.33 & 0.82 & 0.01 & 0.19\\
\midrule
AHU coil & LogReg (stats) & 1.00 & 0.76 & 1.00 & 1.00 & 1.00 & 0.97 & 1.00 & 0.00 & 0.01\\
 & Random forest & 1.00 & 0.66 & 0.98 & 1.00 & 0.81 & 0.86 & 0.89 & 0.05 & 0.01\\
 & Grad.\ boosting & 1.00 & 0.45 & 0.81 & 0.98 & 0.52 & 0.57 & 0.54 & 0.03 & 0.00\\
 & MLP & 1.00 & 0.39 & 0.79 & 0.81 & 0.63 & 0.40 & 0.71 & 0.00 & 0.01\\
 & LSTM & 1.00 & 0.38 & 0.78 & 0.63 & 0.58 & 0.30 & 0.73 & 0.00 & 0.01\\
 & Channel-indep. & 1.00 & 0.59 & 0.93 & 0.96 & 0.76 & 0.81 & 0.95 & 0.00 & 0.00\\
 & Capped CI & 1.00 & 0.47 & 0.96 & 0.98 & 0.41 & 0.31 & 0.50 & 0.01 & 0.06\\
 & Reference rule & 0.97 & 0.37 & 0.74 & 0.78 & 0.38 & 0.74 & 0.90 & 0.01 & 0.01\\
\bottomrule
\end{tabular}}
\end{table}

\begin{table}[h]\centering\small
\caption{Real never-seen faults (LSTM vs.\ reference rule), flag rate by the number of signature sensors the fault moves in the fault direction; one seed.}\label{tab:realov}
\setlength{\tabcolsep}{4pt}
\begin{tabular}{lcccc}
\toprule
Target, detector & 0 of 3 & 1 of 3 & 2 of 3 & 3 of 3\\
\midrule
FCU fan (clean), LSTM & 0.22 & 0.63 & 0.79 & 0.93\\
FCU fan (clean), reference rule & 0.06 & 0.13 & 0.36 & 0.44\\
FCU fan (clean), LSTM + channel near-misses & 0.15 & 0.11 & 0.22 & 0.80\\
FCU valve, LSTM & 0.18 & 0.59 & 0.65 & 1.00\\
FCU valve, reference rule & 0.05 & 0.06 & 0.23 & 0.19\\
AHU coil valve, LSTM & 0.00 & 0.04 & 0.01 & --\\
\bottomrule
\end{tabular}
\end{table}



\section{Reproducibility}
All code (data loaders, splicing, models, theory checks, plotting) is NumPy/scikit-learn based and will be released. The real-data benchmark is public \citep{granderson2023}. Every number in \cref{sec:exp} is produced by the scripts named above from fixed seeds. Real-data numbers come from single runs, reported as produced.
